% !TeX root = ../main.tex

\section{Математическое описание устройства}

\subsection*{3.1 Входной сигнал}
На вход устройства поступает аддитивная смесь гармонического сигнала и шума:
\begin{equation}
    s(t) = \sqrt{P_{s}} \sin(2\pi f_c t) + noise(t)
\end{equation}
где $P_s$ — мощность сигнала, $f_c$ — несущая частота.

\subsection*{3.2 Преобразование в АЦП}
Оцифрованный сигнал представляет собой дискретную последовательность:
\begin{equation}
    x[n] = Q_{na}(s(n T_s))
\end{equation}
где $T_s = 1/f_s$ — шаг дискретизации, $Q_{na}$ — операция нелинейного квантования с разрядностью $n_a$.

\subsection*{3.3 Формирование сигнала ЦГ (NCO)}
С учетом заложенных в математическую модель амплитудных ($amv$, в дБ) и фазовых ($pmv$, в градусах) искажений, синфазная ($I_{NCO}$) и квадратурная ($Q_{NCO}$) составляющие опорного сигнала описываются как:
\begin{equation}
    I_{NCO}[n] = 10^{\frac{amv}{20}} \cos(2\pi f_g n T_s)
\end{equation}
\begin{equation}
    Q_{NCO}[n] = \sin\left(2\pi f_g n T_s + \frac{\pi \cdot pmv}{180}\right)
\end{equation}
Итоговый комплексный сигнал цифрового гетеродина имеет вид:
\begin{equation}
    y_{NCO}[n] = I_{NCO}[n] - j Q_{NCO}[n]
\end{equation}
Вся математика базируется на исходном коде. \textit{При использовании целочисленного формата (\texttt{fixed}) к гетеродину применяется масштабирующий коэффициент усиления $K_{NCO} = 2^{n_g - 1} - 1$.}

\subsection*{3.4 Выходной сигнал (Демодуляция)}
Комплексный демодулированный сигнал на выходе устройства получается путем дискретного перемножения сигнала с АЦП и сигнала гетеродина:
\begin{equation}
    z[n] = x[n] \cdot y_{NCO}[n]
\end{equation}
Данная математическая операция тождественна раздельному формированию двух квадратурных каналов: 
Синфазный канал (I):
\begin{equation}
    Re(z[n]) = x[n] \cdot I_{NCO}[n]
\end{equation}
Квадратурный канал (Q):
\begin{equation}
    Im(z[n]) = - x[n] \cdot Q_{NCO}[n]
\end{equation}

На данном этапе симуляции еще производится построение спектров в разных точках симуляции.
